\documentclass[sigconf]{acmart}

%% NOTE TO AUTHOR: this draft targets the 2nd free5GC World Forum
%% (In-Cooperation with ACM SIGSAC), short-paper track (4 pages,
%% submission and camera-ready). Submissions are NOT double-blind, so
%% author/affiliation fields below should be filled in with real values
%% before compiling for submission. See paper/NOTES.md for what's
%% verified vs. still needs your action (author info, real compilation,
%% CCS codes). The three bibliography entries below were independently
%% verified against their primary sources (arXiv abstract pages, DARPA's
%% own site) before this draft was finalized -- re-check the DOIs/URLs
%% still resolve at submission time, since that's cheap insurance.

\settopmatter{printacmref=true}
\renewcommand\footnotetextcopyrightpermission[1]{}
\acmConference[free5GC World Forum '26]{2nd free5GC World Forum, In-Cooperation with ACM SIGSAC}{December 17--18, 2026}{}

\begin{document}

\title{Reachability-Guided Triage and Upstream-Verified Patching for a Real free5GC Vulnerability}

%% TODO: replace [email@institution.edu] with a real contact email
%% before submission -- not provided, left as a placeholder.
\author{Jothilingam Dheeraj}
\affiliation{%
  \institution{Nanyang Technological University}
  \city{Singapore}
  \country{Singapore}
}
\email{[email@institution.edu]}

\begin{abstract}
Automated vulnerability analysis tools are frequently evaluated against
synthetic reproductions of security bugs rather than the real upstream
codebases those bugs were disclosed in, which limits confidence that
results transfer to practice. We present a case study applying
AST-based reachability filtering and rule-based patch generation to a
real, disclosed vulnerability (CVE-2026-40248, CWE-285) in free5GC's
Unified Data Repository (UDR), a component of a real, standards-compliant
open-source 5G core network. Reachability analysis over the real
$\sim$86KB source file reduces 102 parsed functions to the 4 that are
actually reachable from the vulnerable HTTP entry points, a 96.1\%
reduction in analysis scope. We generate a targeted patch for each
reachable finding and, going beyond isolated-file compilation, verify
the patch against a live clone of the actual \texttt{free5gc/udr}
repository: applying all four generated patches preserves a clean build
against the project's real dependency graph. We report this result
alongside its honest limits: compile verification is not runtime
confirmation, and one of the four findings is only partially corrected
by the rule-based patcher. We discuss what further verification
against the live network function would require.
\end{abstract}

\keywords{5G core network security, vulnerability triage, reachability analysis, program repair, free5GC}

\maketitle
%% Narrow ACM two-column width overflows on a few long inline \texttt{}
%% tokens (package/URL/file names) otherwise -- confirmed by a real
%% compile (4 overfull \hbox warnings), not applied speculatively.
\sloppy

\section{Introduction}

Open-source 5G core network software is increasingly deployed in
research testbeds, private networks, and production trials, yet
security tooling aimed at such codebases is commonly validated against
small synthetic examples rather than the real projects it is meant to
help secure. This gap matters specifically for 5G/6G infrastructure:
network functions such as free5GC's Unified Data Repository (UDR) are
large, real-world Go codebases with external dependencies (HTTP
routing, database drivers, and the project's own internal packages),
and a technique that only ever compiles against an isolated copy of one
file has not demonstrated it survives contact with the actual build.

This paper presents a small, fully reproducible case study that closes
part of that gap for one real, disclosed free5GC vulnerability. We
adopt the reachability-filtered code decomposition methodology of
OpenAnt~\cite{openant}: parse a codebase into functions, build a call
graph, and perform breadth-first traversal from attacker-reachable
entry points. OpenAnt reports this reduces analysis scope by up to
97\% on real projects such as OpenSSL and WordPress; we apply the same
method to a single, real, disclosed free5GC vulnerability. Beyond OpenAnt's own
evaluation, we add a verification step: rather than stopping at
generating a patch, we apply it to a live clone of the actual upstream
repository and confirm the result compiles against the project's real
dependency graph, not a scratch copy of a single file. To our knowledge
this upstream-clone compile verification is not part of OpenAnt's
published evaluation or other prior reachability-based triage work.

Our contributions are: (i) an application of OpenAnt's
reachability-filtered decomposition methodology to a real free5GC
source file against a real, disclosed CVE, with a 96.1\% reduction in
the functions requiring further analysis; (ii)
automatically generated patches for all four reachable findings,
verified to preserve a clean build of the real, live-cloned upstream
project; and (iii) an honest accounting of what this does and does not
establish, including a partially-corrected finding we report rather
than omit.

\section{Background}

\textbf{free5GC} is an open-source, Apache-2.0-licensed implementation
of the 5G core network, built around a Service-Based Architecture
(SBA) in which each network function is an independently versioned
component communicating over defined service-based interfaces. Its
goal is to implement the 5GC defined in 3GPP Release 15 and
beyond~\cite{free5gc}, and it has been a Linux Foundation project since
September 2024, governed by a Technical Steering Committee with major
contributions from National Yang Ming Chiao Tung University. The UDR
(Unified Data Repository) network function exposes the
Nudr\_DataRepository service-based interface, which other network
functions use to read and write subscriber and policy data; this is
reflected in the source layout under \texttt{internal/sbi/}. Because SBA
components are independently versioned, the vulnerability described
below is scoped to the UDR module's own release history, not to a
free5GC project-wide version number.

\textbf{CVE-2026-40248} (GHSA-jgq2-qv8v-5cmj, CVSS 3.1 base score 7.5,
CWE-285 Improper Authorization) affects free5GC/udr versions $\leq$
1.4.2. The handler responsible for creating or updating a Traffic
Influence Subscription checks whether the \texttt{influenceId} path
segment equals the literal string \texttt{"subs-to-notify"} and writes
an HTTP 404 response when it does not, but the function does not
\texttt{return} after writing that response, so execution falls through
and the subscription is created or overwritten regardless of the
check's outcome. An unauthenticated attacker on the 5G Service Based
Interface can therefore create or overwrite arbitrary Traffic Influence
Subscriptions, including attacker-controlled callback URIs, while the
API misleadingly reports 404 Not Found in both the legitimate and the
attack case. A sibling vulnerability in the same file and commit
family, CVE-2026-40246 (GHSA-g9cw-qwhf-24jp, also CWE-285, CVSS 7.5),
shows the identical missing-\texttt{return} pattern in the handler that
deletes a subscription. The real fix for this vulnerability family
(commit \texttt{8668627}\footnote{Full hash:
\texttt{86686276a7e226183ee786e3dd6714ec56c78fda}} on
\texttt{github.com/free5gc/udr}) inserts the missing \texttt{return}
statement across four affected handlers in
\texttt{internal/sbi/api\_datarepository.go}.

\section{Approach}

\subsection{Reachability filtering}

Following OpenAnt's decomposition methodology~\cite{openant}, we parse
the target source with \texttt{tree-sitter}, a real grammar-based
parser rather than a regex or brace-counting heuristic, build a call
graph over all top-level function and method definitions, and perform
a breadth-first traversal from a stated list of entry points: here,
the four real free5GC UDR HTTP handler functions that a
service-based-interface client can invoke directly.
Any function not reachable from those entry points is excluded from
further analysis, on the premise that code an external attacker cannot
reach through the declared interface does not need to be triaged with
the same priority as code that is directly exposed.

\subsection{Triage and patch generation}

Each reachable function is checked against a small library of
rule-based vulnerability patterns. The pattern relevant to this case,
\texttt{find\_missing\_return\_after\_response}, flags a function that
writes an HTTP error response inside a conditional block without a
following \texttt{return} statement in that block, exactly the
control-flow defect described in Section 2. When the pattern matches,
a patch is generated by inserting a \texttt{return} statement
immediately after the flagged response write. This is a targeted,
minimal-diff patch, not a full-function rewrite.

\subsection{Upstream verification}

Prior work verifying a generated patch typically applies it to an
isolated copy of the single affected file and checks that the file
alone still compiles. This does not establish that the patch is
compatible with the project's actual build: a real Go module's
compilation depends on its full dependency graph, declared in
\texttt{go.mod}/\texttt{go.sum}, and on how the patched function
interacts with the rest of the package. We instead clone the real
\texttt{free5gc/udr} repository at the commit immediately preceding the
disclosed fix, confirm the unmodified repository builds
(\texttt{go build ./...}) as a baseline, apply all generated patches
directly to the real \texttt{internal/sbi/api\_datarepository.go} file
inside that clone, and rebuild. This exercises the project's actual
dependency graph, including \texttt{gin-gonic/gin},
\texttt{go.mongodb.org/mongo-driver}, and free5GC's own internal
packages, rather than a synthetic stand-in.

\section{Evaluation}

\textbf{Reachability.} Parsing
\texttt{internal/sbi/\allowbreak api\_datarepository.go} (a real, $\sim$86KB source
file) yields 102 top-level functions. Breadth-first traversal from the
four real UDR entry points reaches exactly 4 of them, a 96.1\%
reduction in the functions a subsequent triage step must consider.

\textbf{Triage.} All 4 reachable functions are correctly flagged by the
missing-\texttt{return} pattern and classified as CWE-285, matching the
real vulnerability class disclosed for this CVE family.

\textbf{Patch generation and upstream verification.} A patch is
generated for all 4 flagged functions. The unmodified clone builds
cleanly as a baseline. With all 4 generated patches applied directly to
the real, live-cloned repository, \texttt{go build ./...} again
succeeds against the project's actual dependency graph. As far as we
are aware, this is the first time a patch produced by this
reachability pipeline has been verified against a real upstream
module's full dependency graph rather than an isolated single-file
copy.

\textbf{A limitation surfaced by measurement, not assumption.} One of
the four flagged functions requires \emph{two} separate
\texttt{return} insertions in the real, disclosed fix (two
independent conditional blocks in the same function body), while our
rule-based patcher returns only the first match per function. That
function is therefore only partially corrected: the build still
succeeds (a partial fix does not break compilation), but the
vulnerability is not fully closed for that one handler. We report this
rather than omit it, since it is a genuine, measured limitation of a
single-finding-per-pattern-match design, not a flaw discovered after
publication.

\section{Discussion and Limitations}

A clean build after patching establishes that a generated patch is
syntactically valid and compatible with the real project's dependency
graph. It does not establish that the patch is correct at runtime.
Confirming that the vulnerability is actually closed would require
running the real free5GC UDR service and sending it a live HTTP
request, which in turn requires the rest of the deployment this network
function depends on: a MongoDB backend for its data store and
registration with the free5GC NRF for service discovery. Standing up
that deployment was outside the scope of this case study and is the
most direct next step for this work. We also note this is a
single-project, single-vulnerability-family case study (n=1); we do not
claim the reachability reduction or patch-generation success rate
observed here generalizes to free5GC broadly or to other 5G core
implementations without further evaluation.

\section{Related Work}

Our reachability filtering directly follows OpenAnt~\cite{openant},
which combines reachability-based code decomposition with LLM-driven
adversarial verification and sandboxed dynamic testing to discover
previously unknown vulnerabilities across projects including OpenSSL,
WordPress, and Flowise, reporting up to 97\% reduction in analysis
scope. We apply the same decomposition step to a single, real,
disclosed CVE rather than searching for unknown bugs, and add a
verification step OpenAnt's own evaluation does not report: confirming
a generated patch against a live clone of the real upstream repository,
not a sandboxed exploit environment discarded after use.

DARPA's AI Cyber Challenge (AIxCC) produced seven open-sourced cyber
reasoning systems combining fuzzing, static analysis, and LLM agents
for autonomous vulnerability discovery and patching~\cite{aixcc}.
OSS-CRS extends this line of work by porting AIxCC-era systems to real
OSS-Fuzz projects under resource budgets~\cite{osscrs}. Our work is
narrower in scope than either: we do not attempt discovery of unknown
vulnerabilities, only reachability-filtered triage and patching of an
already-disclosed one. It shares their emphasis on evaluating
against real, not synthetic, targets.

\section{Conclusion}

We presented a case study demonstrating that reachability-filtered
static analysis and automated patch generation can be applied to a
real, disclosed free5GC vulnerability with verification against the
project's actual upstream repository, not a synthetic reproduction.
This reduced 102 real functions to the 4 that matter and confirmed
generated patches for all 4 preserve a clean build of the live cloned
project. We report this alongside its real limits: compile-time
verification is not runtime confirmation, and one finding is only
partially corrected by the current rule-based patcher. We see this as
a concrete, if narrow, step toward evaluating automated vulnerability
tooling against real 5G core network software rather than stand-ins
built for the purpose.

\begin{thebibliography}{9}

\bibitem{openant}
Nahum Korda and Gadi Evron. 2026. OpenAnt: LLM-Powered Vulnerability
Discovery Through Code Decomposition, Adversarial Verification, and
Dynamic Testing. \textit{arXiv preprint arXiv:2606.19149} (2026).
\url{https://arxiv.org/abs/2606.19149}

\bibitem{free5gc}
free5GC Project. \url{https://free5gc.org/}.

\bibitem{aixcc}
DARPA. 2025. AI Cyber Challenge marks pivotal inflection point for
cyber defense. (August 2025).
\url{https://www.darpa.mil/news/2025/aixcc-results}

\bibitem{osscrs}
Andrew Chin, Dongkwan Kim, Yu-Fu Fu, Fabian Fleischer, Youngjoon Kim,
HyungSeok Han, Cen Zhang, Brian Junekyu Lee, Hanqing Zhao, and Taesoo
Kim. 2026. OSS-CRS: Liberating AIxCC Cyber Reasoning Systems for
Real-World Open-Source Security. \textit{arXiv preprint
arXiv:2603.08566} (2026). \url{https://arxiv.org/abs/2603.08566}

\end{thebibliography}

\end{document}
